\subsection{So sánh đánh giá retrieval rời rạc và đánh giá end-to-end}

Trong đồ án này, kết quả thực nghiệm được đọc từ hai nhóm artifact chính:
\texttt{evaluation/results/custom\_eval} và
\texttt{evaluation/results/e2e\_custom\_eval}. Hai nhóm này không đo cùng một
đối tượng chất lượng, vì vậy cần tách rõ khi diễn giải.

Nhóm \texttt{custom\_eval} dùng script
\texttt{evaluate\_retrieval\_custom.py} để đánh giá riêng tầng truy hồi. Với mỗi
câu hỏi, script chạy router classifier, chọn collection bằng
\texttt{CollectionSelector}, lấy ứng viên bằng hybrid search
Qdrant--Elasticsearch, rerank bằng cross-encoder và tính Hit, Precision,
Recall, MRR, nDCG dựa trên tập \texttt{evidence\_chunk\_ids} trong dataset.
Nhóm này không sinh câu trả lời cuối cùng, do đó không có các metric về
faithfulness, completeness hay mức khớp với đáp án vàng.

Ngược lại, nhóm \texttt{e2e\_custom\_eval} dùng script
\texttt{evaluate\_e2e\_pipeline.py} để gọi trực tiếp
\texttt{RAGPipeline.query\_v3()}. Đây là entrypoint của luồng hỏi đáp
end-to-end, bao gồm routing, reflection, chọn collection, hybrid retrieval,
reranking, lọc hiệu lực văn bản, mở rộng parent context và generation. Sau khi
có câu trả lời, evaluator vừa tính lại các metric retrieval trên danh sách
source cuối cùng, vừa dùng \texttt{SelfEvaluator} để đánh giá relevance,
faithfulness và completeness, đồng thời dùng một LLM judge để so sánh câu trả
lời sinh ra với \texttt{gold\_answer}.

\begin{table}[H]
\centering
\small
\begin{tabular}{|p{3.2cm}|p{4.8cm}|p{6.1cm}|}
\hline
\textbf{Nhóm đánh giá} & \textbf{Đối tượng đo} & \textbf{Ý nghĩa chính} \\
\hline
\texttt{custom\_eval} & Retrieval rời rạc, chưa sinh câu trả lời. & Đo khả năng đưa đúng chunk vàng vào top-K; phù hợp để chẩn đoán router, hybrid search và reranker. \\
\hline
\texttt{e2e\_custom\_eval} & Toàn bộ pipeline RAG có sinh câu trả lời. & Đo chất lượng trả lời cuối cùng, gồm mức bám context, độ đầy đủ và mức khớp đáp án vàng; retrieval chỉ là một phần trong luồng này. \\
\hline
\end{tabular}
\caption{Khác biệt giữa hai nhóm artifact đánh giá}
\end{table}

Với dataset \texttt{ITE6\_rag\_evaluation\_dataset\_no\_parent\_evidence}, hai
nhóm artifact có cùng 26 câu hỏi nên có thể đối chiếu trực tiếp. Bảng dưới cho
thấy chất lượng retrieval của hai cách chạy tương đối gần nhau ở Hit@5, trong
khi run end-to-end có thêm các chỉ số chất lượng câu trả lời.

\begin{table}[H]
\centering
\small
\begin{tabular}{|p{4.9cm}|r|r|r|r|r|}
\hline
\textbf{Run} & \textbf{Hit@5} & \textbf{Recall@5} & \textbf{MRR@5} & \textbf{nDCG@5} & \textbf{Latency TB} \\
\hline
\texttt{custom\_eval/ITE6} & 73,08\% & 68,59\% & 62,31\% & 62,76\% & 4.375,3 ms \\
\hline
\texttt{e2e\_custom\_eval/ITE6\_\_min\_top\_k\_7\_hyde} & 73,08\% & 71,15\% & 54,81\% & 57,95\% & 10.701,2 ms \\
\hline
\end{tabular}
\caption{So sánh retrieval trên cùng dataset ITE6}
\end{table}

Ở run end-to-end của ITE6, mặc dù MRR@5 và nDCG@5 thấp hơn run retrieval rời
rạc, chất lượng câu trả lời vẫn ở mức tốt: answer relevance đạt 100,00\%,
faithfulness đạt 84,62\%, completeness đạt 88,46\% và tỷ lệ khớp đúng với đáp
án tham chiếu đạt 80,77\%. Điều này cho thấy thứ hạng chunk vàng theo
\texttt{evidence\_chunk\_ids} không phải lúc nào cũng phản ánh đầy đủ chất
lượng câu trả lời cuối cùng.

Nguyên nhân chính nằm ở cách xây dựng ground truth hiện tại. Dataset đánh giá
được sinh theo từng file chunk trong \texttt{data/}: với mỗi file chunk, hệ
thống sinh ra một file dataset tương ứng và gán \texttt{evidence\_chunk\_ids}
từ chính file nguồn đó. Tuy nhiên, trong kho tri thức thực tế, cùng một thông
tin học vụ có thể xuất hiện ở nhiều file khác nhau, hoặc được diễn đạt lại
trong các văn bản liên quan thuộc cùng collection hay collection khác. Vì vậy,
một câu trả lời vẫn có thể được grounding tốt từ context đúng nhưng context đó
không trùng với danh sách chunk vàng đã gán trong dataset. Khi đó Hit@K,
Recall@K, MRR@K hoặc nDCG@K có thể bị đánh thấp, dù câu trả lời thực tế vẫn
faithful và đúng với nội dung cần trả lời.

Ngoài ra, trong luồng end-to-end, parent context được dùng để làm giàu ngữ cảnh
cho LLM sau bước rerank. Theo cài đặt trong
\texttt{evaluate\_e2e\_pipeline.py}, parent expansion không chèn thêm chunk
parent vào danh sách ranked chunks và không thay đổi thứ tự kết quả truy hồi.
Do đó, parent context có thể cải thiện grounding và completeness của câu trả
lời, nhưng không trực tiếp làm tăng các metric dựa trên so khớp
\texttt{evidence\_chunk\_ids}. Đây là một giới hạn quan trọng khi diễn giải
retrieval metrics trong hệ thống RAG có mở rộng ngữ cảnh.

\subsection{Kết quả tổng hợp được sử dụng trong báo cáo}

Để bảng tổng hợp phản ánh cấu hình ổn định của hệ thống, báo cáo chính chỉ sử
dụng các run có đủ artifact \texttt{summary.json}, được chạy với cấu hình
\texttt{min\_top\_k\_7}, và loại khỏi bảng chính các lô có chất lượng generation
thấp hoặc không đại diện cho cấu hình cuối. Cụ thể, nhóm tổng hợp chọn các run
có faithfulness tối thiểu 80\% và tỷ lệ \texttt{ref\_match=correct} tối thiểu
75\%. Các artifact còn lại vẫn có giá trị để phân tích lỗi, nhưng không được
dùng làm số liệu đại diện trong bảng kết quả chính.

Với tiêu chí trên, nhóm end-to-end được báo cáo gồm 27 run, tổng cộng 876 câu
hỏi. Kết quả retrieval đạt Hit@5 = 89,16\%, Recall@5 = 86,64\%, MRR@5 =
78,10\% và nDCG@5 = 79,07\%. Quan trọng hơn, các metric chất lượng câu trả lời
đạt mức cao: relevance = 99,77\%, faithfulness = 89,61\%, completeness =
94,29\%, hallucination rate = 6,62\%, và tỷ lệ trả lời đúng hoàn toàn so với
đáp án vàng đạt 87,79\%.

\begin{table}[H]
\centering
\small
\begin{tabular}{|p{5.1cm}|r|}
\hline
\textbf{Metric end-to-end chọn lọc} & \textbf{Giá trị} \\
\hline
Số run / số câu hỏi & 27 run / 876 câu \\
\hline
Hit@5 / Recall@5 & 89,16\% / 86,64\% \\
\hline
MRR@5 / nDCG@5 & 78,10\% / 79,07\% \\
\hline
Answer relevance & 99,77\% \\
\hline
Faithfulness & 89,61\% \\
\hline
Completeness & 94,29\% \\
\hline
Hallucination rate & 6,62\% \\
\hline
Ref match correct & 87,79\% \\
\hline
\end{tabular}
\caption{Kết quả end-to-end trên nhóm artifact ổn định dùng cho báo cáo chính}
\end{table}

Từ các kết quả trên, có thể kết luận rằng các metric retrieval nên được xem là
metric chẩn đoán chất lượng truy hồi theo chunk-id, không nên xem là thước đo
duy nhất của chất lượng hệ thống. Trong bối cảnh dữ liệu học vụ có nhiều văn
bản trùng lặp hoặc bổ sung lẫn nhau, faithfulness, completeness và reference
match phản ánh sát hơn chất lượng người dùng nhận được ở đầu ra cuối cùng. Đây
cũng là lý do báo cáo ưu tiên đánh giá cao các metric end-to-end, đặc biệt là
faithfulness và mức khớp đáp án vàng, đồng thời ghi nhận giới hạn của bộ
ground-truth retrieval hiện tại như một điểm cần cải thiện trong các vòng đánh
giá sau.
